% ============================================================
%  Příprava na maturitu – Mechatronika
%  Kompiluj: pdflatex maturita_mechatronika.tex
%  Kódování: UTF-8
% ============================================================

\documentclass[a4paper,10pt]{article}

% --- Balíčky ---
\usepackage[T1]{fontenc}
\usepackage[utf8]{inputenc}
\usepackage[czech]{babel}
\usepackage[left=2.4cm, right=1.8cm, top=1.8cm, bottom=1.8cm]{geometry}
\usepackage{lmodern}
\usepackage{microtype}
\usepackage{parskip}
\usepackage{titlesec}
\usepackage{enumitem}
\usepackage{xcolor}
\usepackage{hyperref}
\usepackage{tabularx}
\usepackage{booktabs}
\usepackage{array}
\usepackage{tikz}
\usepackage{eso-pic}
\usepackage{mdframed}
\usepackage{amsmath}
\usepackage{amssymb}
\usepackage{pgfplots}
\pgfplotsset{compat=1.18}
\usetikzlibrary{arrows.meta, positioning, shapes.geometric, calc}

% --- Barvy ---
\definecolor{accent}{HTML}{03254E}
\definecolor{stripeblue}{HTML}{568EA3}
\definecolor{medgray}{HTML}{888888}
\definecolor{lightblue}{HTML}{EAF2F8}
\definecolor{lightyellow}{HTML}{FDFBE4}
\definecolor{lightgreen}{HTML}{EAF7EA}

% --- Modrý pruh vlevo ---
\AddToShipoutPictureBG{%
  \begin{tikzpicture}[remember picture, overlay]
    \fill[stripeblue]
      (current page.north west)
      rectangle
      ([xshift=10mm] current page.south west);
    \fill[accent!60]
      ([xshift=10mm] current page.north west)
      rectangle
      ([xshift=12mm] current page.south west);
  \end{tikzpicture}%
}

\hypersetup{colorlinks=true, urlcolor=accent, linkcolor=accent}

\titleformat{\section}
  {\large\bfseries\color{accent}}
  {\thesection.}
  {0.5em}
  {}
  [\color{accent}\titlerule]

\titlespacing{\section}{0pt}{16pt}{6pt}

\newmdenv[
  backgroundcolor=lightblue,
  linecolor=stripeblue,
  linewidth=1pt,
  roundcorner=4pt,
  innerleftmargin=8pt,
  innerrightmargin=8pt,
  innertopmargin=6pt,
  innerbottommargin=6pt,
  skipabove=4pt,
  skipbelow=4pt
]{pojmybox}

\newmdenv[
  backgroundcolor=lightyellow,
  linecolor=accent!40,
  linewidth=1pt,
  roundcorner=4pt,
  innerleftmargin=8pt,
  innerrightmargin=8pt,
  innertopmargin=6pt,
  innerbottommargin=6pt,
  skipabove=4pt,
  skipbelow=8pt
]{vysvetlenibox}

\newmdenv[
  backgroundcolor=lightgreen,
  linecolor=accent!30,
  linewidth=1pt,
  roundcorner=4pt,
  innerleftmargin=8pt,
  innerrightmargin=8pt,
  innertopmargin=6pt,
  innerbottommargin=6pt,
  skipabove=4pt,
  skipbelow=8pt
]{vzorcebox}

\pagestyle{plain}

% ============================================================
\begin{document}

% --- Záhlaví ---
\begin{minipage}[t]{0.75\textwidth}
    {\Huge\bfseries\color{accent} Příprava na maturitu}\\[4pt]
    {\large\color{stripeblue} Mechatronika}\\[2pt]
    \textcolor{medgray}{\small Straka Jakub, Suchánek Lukáš \quad SPŠ Prosek \quad 2026}
\end{minipage}
\hfill
\begin{minipage}[t]{0.22\textwidth}
    \raggedleft
    \begin{tikzpicture}
        \draw[color=stripeblue, rounded corners=4pt, line width=1pt]
            (0,0) rectangle (3,1.2)
            node[midway, align=center, color=accent, font=\small\bfseries]
            {Otázek: 22};
    \end{tikzpicture}
\end{minipage}

\vspace{6pt}
\noindent\color{accent}\rule{\linewidth}{1.5pt}
\vspace{2pt}

% ============================================================
\section{Mechanismy obecného pohybu a tvrdá automatizace}

\begin{pojmybox}
\textbf{Klíčové pojmy:}
vačka, páka, kliková hřídel, kulisový mechanismus, Genevský mechanismus,
tvrdá automatizace, přenosový stroj, karuselový stroj, stupeň volnosti,
Grüblerova podmínka
\end{pojmybox}

\begin{vysvetlenibox}
\textbf{Stupně volnosti mechanismu}

\begin{vzorcebox}
Grüblerova podmínka (rovinný mechanismus):
$$F = 3(n-1) - 2p_1 - p_2$$
kde $n$ = počet členů (včetně rámu), $p_1$ = počet jednoduchých kloubů,
$p_2$ = počet dvojkloubů. Mechanismus: $F=1$, struktura: $F=0$.
\end{vzorcebox}

\textbf{Základní mechanismy}

\begin{center}
\begin{tabularx}{0.95\linewidth}{@{} l X l @{}}
    \toprule
    \textbf{Mechanismus} & \textbf{Funkce} & \textbf{Aplikace} \\
    \midrule
    Vačkový      & Rotace $\to$ libovolný pohyb sledovače;
                   profil vačky určuje zákon pohybu & Ventily motoru, podavače \\
    Klikový       & Rotace $\to$ přímočarý pohyb & Pístové motory, lisy \\
    Genevský      & Přerušovaný pohyb (1 otočení $\to$ 1 krok)
                  & Indexovací stoly, filmové projektory \\
    Páka          & Přenos a změna síly/pohybu & Záchyty, spínače \\
    Kulisový      & Kluzná i rotační vazba & Variátory rychlosti \\
    \bottomrule
\end{tabularx}
\end{center}

\textbf{Tvrdá automatizace -- charakteristika}\\
Stroje pro hromadnou výrobu s neměnným sledem operací.
Maximální produktivita, ale nulová flexibilita při změně výrobku.

\begin{center}
\begin{tabularx}{0.95\linewidth}{@{} l X @{}}
    \toprule
    \textbf{Typ stroje} & \textbf{Popis} \\
    \midrule
    Přenosový (in-line) & Obrobky se posunují lineárně od stanice ke stanici \\
    Karuselový (rotační) & Otočný stůl, obrobky se otáčejí kolem osy \\
    Gnezdový (cellular)  & Obrobky se zpracovávají na jednom místě \\
    \bottomrule
\end{tabularx}
\end{center}
\end{vysvetlenibox}

% ============================================================

% ============================================================
\end{document}
